%==============================================================================
% Chapitre 22 - Systèmes de conduite de tir
%==============================================================================

\section{Introduction}

Un système de conduite de tir (SCT) est un ensemble de moyens destinés à
assister le tireur dans l'engagement d'une cible.

Son objectif est de déterminer la solution de tir la plus appropriée à partir
des informations disponibles.

Les systèmes modernes combinent généralement :

\begin{itemize}
    \item des capteurs ;
    \item des calculateurs ;
    \item des systèmes optiques ;
    \item des interfaces homme-machine ;
    \item des dispositifs de pointage.
\end{itemize}

\begin{aretenir}

Un système de conduite de tir transforme des informations de mesure en
corrections de tir exploitables.

\end{aretenir}

\section{Évolution historique}

Les premiers systèmes de conduite de tir reposaient essentiellement sur des
tables de tir et des calculs manuels.

Avec le développement de l'électronique sont apparus :

\begin{itemize}
    \item les calculateurs analogiques ;
    \item les calculateurs numériques ;
    \item les systèmes automatisés ;
    \item les systèmes intégrés modernes.
\end{itemize}

Aujourd'hui, certains systèmes sont capables de calculer une solution de tir
en quelques fractions de seconde.

\section{Fonctions principales}

Un système de conduite de tir doit généralement :

\begin{enumerate}
    \item acquérir les données nécessaires ;
    \item calculer une solution balistique ;
    \item assister le pointage ;
    \item suivre la cible ;
    \item enregistrer certaines informations.
\end{enumerate}

Selon les applications, certaines fonctions peuvent être absentes ou
automatisées.

\section{Les données d'entrée}

Le calcul d'une solution de tir nécessite de nombreuses informations.

Parmi les plus courantes :

\begin{itemize}
    \item distance à la cible ;
    \item vitesse de la cible ;
    \item direction de déplacement ;
    \item vitesse initiale ;
    \item conditions atmosphériques ;
    \item caractéristiques du projectile ;
    \item position géographique.
\end{itemize}

\section{Mesure de la distance}

La distance constitue souvent l'information la plus importante.

Elle peut être obtenue par :

\begin{itemize}
    \item estimation visuelle ;
    \item télémètre optique ;
    \item télémètre laser ;
    \item systèmes radar.
\end{itemize}

Une erreur sur la distance peut entraîner une erreur importante sur la
solution de tir.

\section{Mesure des conditions atmosphériques}

Les systèmes modernes peuvent intégrer :

\begin{itemize}
    \item température ;
    \item pression ;
    \item humidité ;
    \item vent.
\end{itemize}

Ces informations permettent d'améliorer la précision des calculs
balistiques.

\section{Capteurs d'attitude}

La plateforme peut être inclinée ou en mouvement.

Des capteurs spécifiques permettent de mesurer :

\begin{itemize}
    \item roulis ;
    \item tangage ;
    \item lacet ;
    \item accélérations.
\end{itemize}

Ces informations sont essentielles pour les plateformes mobiles.

\section{Le calculateur balistique}

Le calculateur constitue le cœur du système.

Son rôle consiste à déterminer :

\begin{itemize}
    \item la correction en site ;
    \item la correction en dérive ;
    \item l'avance sur cible mobile ;
    \item diverses corrections complémentaires.
\end{itemize}

Le niveau de complexité dépend fortement de l'application considérée.

\section{Solution de tir}

La solution de tir correspond à l'ensemble des corrections nécessaires pour
atteindre la cible.

Elle peut inclure :

\begin{itemize}
    \item la chute du projectile ;
    \item l'effet du vent ;
    \item le mouvement de la cible ;
    \item le mouvement de la plateforme ;
    \item les effets géophysiques.
\end{itemize}

\section{Cibles mobiles}

L'engagement d'une cible mobile nécessite l'introduction d'une avance.

Cette avance dépend notamment :

\begin{itemize}
    \item du temps de vol ;
    \item de la vitesse de la cible ;
    \item de sa direction de déplacement.
\end{itemize}

\begin{exemple}

Une cible se déplaçant perpendiculairement à la ligne de tir nécessite une
avance plus importante qu'une cible se déplaçant directement vers le tireur.

\end{exemple}

\section{Suivi automatique}

Les systèmes avancés peuvent assurer le suivi automatique d'une cible.

Ils utilisent généralement :

\begin{itemize}
    \item des algorithmes de poursuite ;
    \item des caméras ;
    \item des capteurs optroniques ;
    \item des radars.
\end{itemize}

Ces fonctions sont particulièrement répandues sur les plateformes modernes.

\section{Stabilisation et conduite de tir}

Le système de conduite de tir travaille souvent en liaison avec le système
de stabilisation.

Les deux systèmes échangent des informations afin de :

\begin{itemize}
    \item maintenir la visée ;
    \item compenser les mouvements ;
    \item améliorer la précision.
\end{itemize}

\section{Architecture moderne}

Un système moderne comprend souvent :

\begin{itemize}
    \item une centrale inertielle ;
    \item un télémètre laser ;
    \item plusieurs capteurs météorologiques ;
    \item un calculateur balistique ;
    \item des viseurs optroniques ;
    \item des interfaces utilisateur.
\end{itemize}

L'ensemble fonctionne comme un système intégré.

\section{Conduite de tir et simulation}

Les systèmes de simulation reproduisent fréquemment tout ou partie d'un SCT.

Ils peuvent modéliser :

\begin{itemize}
    \item les capteurs ;
    \item les erreurs de mesure ;
    \item les algorithmes ;
    \item les délais de calcul ;
    \item les interfaces opérateur.
\end{itemize}

Cette modélisation est essentielle pour les simulateurs d'entraînement.

\begin{simulation}

Dans certains simulateurs, le moteur balistique et le système de conduite de
tir sont séparés.

Le SCT fournit une solution de tir tandis que le moteur balistique calcule
la trajectoire réelle du projectile.

\end{simulation}

\section{Limites des systèmes de conduite de tir}

Même les systèmes les plus avancés présentent des limites.

Les erreurs peuvent provenir :

\begin{itemize}
    \item des mesures ;
    \item des modèles utilisés ;
    \item des capteurs ;
    \item des hypothèses de calcul.
\end{itemize}

La qualité d'une solution dépend directement de la qualité des données
d'entrée.

\section{À retenir}

\begin{aretenir}

\begin{itemize}
    \item Un SCT calcule une solution de tir à partir de données mesurées.
    \item Le calculateur balistique constitue le cœur du système.
    \item Les capteurs fournissent les informations nécessaires au calcul.
    \item Les cibles mobiles nécessitent une avance de tir.
    \item La qualité d'une solution dépend de la qualité des mesures.
\end{itemize}

\end{aretenir}
